\section{Technical Overview}
\label{sec:technical-overview}

\subsection{The ledger}
\label{sec:overview-ledger}

We use the UC framework to formulate our ledger assumption and security
goal.  Here we sketch the ledger functionality $\FLedger$; its complete
interface is deferred to the formal treatment.

\begin{description}
  \item[A synchronous public record]
    The ledger maintains an append-only public list $L$, initially containing
    a common reference string $\crs$ sampled from a specified setup
    distribution $\Setup$.  All parties see the same list at the same time.
    An honest posting is atomic: the adversary cannot inspect it and then
    prevent its inclusion, withdraw it on corruption, or make another
    message overtake it.  The setup distribution is a protocol parameter,
    not a particular cryptographic construction built into the ledger.

  \item[Growing registration and corruption]
    Parties may register throughout the execution.  The ledger generates
    their public keys, records them in $L$, and retains the corresponding
    secret material and random tapes as protected state.  A registered
    tuple contains a delivery-encryption key $\ncek$ and a signature
    verification key $\sigvk$; their use is determined by the protocol
    role.  Key generation is honest even for a party selected for initial
    corruption.  Each newly
    registered computation party is independently corrupt with probability
    $p<1/3$; if so, its state is disclosed to the adversary.  Application
    parties---the input providers and output recipients---are not subject
    to this initial random corruption experiment.

    Registration may take place in large batches long before the keys
    are used.  A new protocol role requires an unused party, not a newly
    registered key.

    An initially honest computation party or input provider is protected
    from targeted corruption until it posts a message.  Once it posts,
    the adversary may corrupt it adaptively and obtain its retained state.
    Output recipients are instead subject to chosen adaptive corruption
    at any time, whether or not they have posted or received their output.
    Chosen corruptions are unbounded; $p$ governs only initial computation
    corruptions.  Registration does not count as speaking.

  \item[One opportunity to speak]
    A fixed, protocol-dependent schedule determines from $L$ which parties
    may post in each synchronous communication round.  They share a
    deadline; there is no separate slot for each party.  Each party may
    post at most once.  Silence by the deadline produces a $\NoMsg$
    record, allowing the execution to continue with the protocol's
    prescribed default.  This ledger-generated record is not a
    message from the party and does not make it eligible for post-message
    corruption.  A corrupt party may post any message;
    interpreting messages and ignoring malformed ones is the protocol's
    responsibility, not the ledger's.

  \item[Private computation without erasures] 

    Write $\sk_P$ for the secret key material retained for party $P$.
    An honest party may privately submit an efficient algorithm $Y$ and
    receive the result
    \[
      y\gets Y(\sk_P,L).
    \]
    These private evaluations may be repeated with different algorithms;
    they neither count as speaking nor change the party's corruption status.

    During its single posting opportunity, the party may instead privately
    submit an efficient algorithm $M$, with any private input embedded in
    its description.  The ledger computes
    \[
      m\gets M(\sk_P,L)
    \]
    and appends $m$ if $m\neq\bot$; otherwise the party remains silent.
    For parties protected until speaking, that protection ends only after
    the message is irrevocably on the ledger.  No state is erased:
    corruption reveals the
    secret material, key-generation tape, and retained private-call history,
    including the submitted algorithms and their computation randomness.
\end{description}

These are deliberately idealized assumptions.  Initial random computation-party
corruption models hidden random assignment of roles to a population with a
fixed compromised fraction; speaking makes a previously protected party available
for targeted compromise.  The model does not allow the adversary to move
its corruption budget between still-silent computation parties.
Knowing a party's public key or its assigned role, including an input role,
does not by itself reveal where its physical machine can be found.  A larger
context might reveal that location and should then permit targeted corruption
even before speaking.  Capturing such contexts requires separating public
keys from physical machines; this composition issue is outside the present
abstraction.

Keeping secret state inside $\FLedger$ is UC bookkeeping, not a proposal
that blockchains hold users' secrets.  It avoids returning secret keys
as intermediate interface outputs visible to the environment.
Although $Y(\sk_P,L)=\sk_P$ would disclose the key, the protocol specifies
honest parties' algorithms.  Garblings, committees, signature checks,
and ciphertext interpretation belong to that protocol, not the ledger.

\subsection{The ideal functionality}
\label{sec:overview-functionality}

Our goal is to realize an ideal functionality $\FMPC$ for ongoing
computation in synchronous I/O rounds.  It holds a secret state $\sigma$,
initially the empty string.  Its interface concerns only input providers
and output recipients: the computation parties and their initial random
corruptions are implementation details and do not appear in this ideal task.

\begin{description}
  \item[The round function]
    In each round, the parties specify a public Boolean circuit $F$ that
    maps the current state and new input bits to a new state and output
    bits, with these groups of wires explicitly designated.  We
    restrict our security claim to executions in which the honest parties
    specify the same $F$; agreement on $F$ is not a task our protocol solves.
    Circuits may be chosen as the computation proceeds.

  \item[Input groups and output wires]
    The circuit's external input wires are partitioned into groups.
    Group $a$ has a public length $\ell_{s,a}$ in round $s$ and a
    designated provider $\pid_a$, which supplies that entire input string
    in one message.  Each output wire $i$ still has its own designated
    recipient.  The functionality notifies parties of their assignments;
    these assignments and lengths are also visible to the environment
    and adversary.  The state wires have no assigned application party.

    In the implementation, an application identifier is its registered
    public-key tuple $\pid=(\ncek,\sigvk)$.  Input providers use the signing
    component and output recipients the delivery-encryption component;
    corruption exposes the retained secret material and tapes for both.
    Each input group and output wire uses its own previously unused
    party, even when several roles serve the same application user.

  \item[Inputs]
    Provider $a$ supplies a private string
    $x_a\in\bits^{\ell_{s,a}}$.  Honest inputs are supplied at the
    common start of the I/O round.  During the public input-posting phase,
    the functionality notifies the adversary which providers supplied
    an input, without revealing the strings; the implementation likewise
    reveals whether a provider posts.  The first correctly sized input
    is accepted, and subsequent submissions cannot change it.
    Malformed submissions are ignored.  If no input is supplied by the
    deadline, set $x_a=0^{\ell_{s,a}}$.
    The concatenation of these strings is the circuit input $x$.

  \item[State update and outputs]
    Once the inputs are fixed, the functionality computes
    \[
      (\sigma',y)=F(\sigma,x).
    \]
    It retains $\sigma'$ and fixes each output bit for its assigned
    recipient; private delivery follows in the output phase.  Neither
    the state nor honest parties' inputs and outputs are made public.

  \item[Adaptive corruption]
    Corruptions obey the admissibility rule in the ledger sketch above.
    In particular, the adversary may choose to corrupt any output
    recipient at any time, including one that never speaks.
    An allowed corruption exposes the affected party's retained inputs
    and its already-computed outputs, including outputs pending private
    delivery.  It does not itself disclose the functionality's internal
    state or another party's private information.  If a recipient is
    already corrupt when its output is computed, that output is given to
    the adversary.  An input already accepted
    cannot be replaced by corrupting its provider later.  We defer the
    precise scheduling and corruption bookkeeping to the formal definition.
\end{description}

\subsection{Gadgets}
\label{sec:overview-gadgets}

Each I/O round uses a fresh garbled circuit implementing $F$ together
with input decryption, state transfer, and output delivery.  We describe
the components before assembling the protocol.

\begin{description}
  \item[Public sampling without a private garbler]
    For our protocol, $\crs$ contains universal-sampler parameters $U$.
    A public protocol rule determines an exact distribution description
    $d$ from the round circuit, input lengths and assignments, registered
    signing and delivery keys, committee rosters, and a nonce $Q$.
    All fields other than $Q$ are fixed on the ledger before its release.
    Everyone computes the same $\Sample(U,d)$: a garbled circuit and
    encrypted information for its participants.  The distribution returns
    neither its random tape nor the garbler's encoding key, so no
    participating machine holds all the generated secrets.  After a
    nonce-only bootstrap, we use one main sample per I/O round.

    Including delivery keys in $d$ is essential.  Replacing one of them
    changes the description and, in the universal sampler's ideal
    experiment, gives a sample generated with independent randomness.
    Honest parties supply inputs only to the instance specified by
    this rule from the ledger.  The adversary may generate and evaluate other
    instances, but does not thereby obtain honest inputs to those instances.

  \item[A shared coin for the nonce]
    Each round $s$ has a nonce $Q_s$, chosen uniformly from $\bits^\lambda$
    and held as a robust secret sharing by a dedicated committee.
    Its members have no other role and receive encrypted share packets,
    just as for label committees below.  They speak only after all other
    fields of $d_s$ have been fixed.

    In one synchronous release round, honest members publish their
    decryption keys.  At the deadline, everyone reconstructs $Q_s$ from
    the recovered shares.  Within the sharing threshold, bad or missing
    contributions cannot change the result or prevent reconstruction.
    In particular, corrupt members cannot steer the nonce by choosing
    their contributions after seeing honest ones.

    The description $d_s$ is a fixed encoding of the agreed fields and
    $Q_s$.  Until release, the hidden nonce protects it against advance
    sampling queries.  The adversary's other queries remain unrestricted.

    The sample for round $s$ draws an independent $Q_{s+1}$ and encrypts
    its shares to a committee reserved in $d_s$.  These packets are
    available without evaluating the garbling or knowing the next
    application circuit and recipients.  A nonce-only bootstrap supplies
    shares of $Q_1$.  Thus no speaking dealer holds the entire nonce,
    and all committees may use previously registered keys.

  \item[Signed encrypted inputs]
    For each input provider $a$ in round $s$, the sampling procedure
    generates a fresh SNCE key pair $(\inputpk_{s,a},\inputsk_{s,a})$.
    It publishes $\inputpk_{s,a}$ and embeds $\inputsk_{s,a}$ only
    inside the garbled circuit.  This pair serves the provider's entire
    input string and belongs to the sampled computation, not a registered
    party.  It is distinct from signing, delivery, and state-transfer keys.

    Once its public key is available, provider $a$ encrypts its input:
    \[
      c^{\mathsf{in}}_{s,a}\sample\Enc(\inputpk_{s,a},x_a).
    \]
    It signs an unambiguous encoding
    $\InputRecord(\sid,s,a,d_s,c^{\mathsf{in}}_{s,a})$ and posts
    the ciphertext and signature together.  Here $\sid$ identifies the
    protocol session, $s$ the round, $a$ the input group, and $d_s$ the
    exact sample description.  All parties verify the signature under
    the provider's registered $\sigvk$; signatures are not themselves
    garbled inputs.

    The first correctly sized ciphertext with a valid signature and the
    right context, received before the common input deadline, is accepted
    irrevocably; other records are ignored.  Later signing-key disclosure
    cannot replace that ciphertext or reopen the deadline, so an ordinary
    one-time signature suffices.

    The circuit decrypts each accepted ciphertext internally to obtain
    $x_a$.  If the chosen encryption has randomized decryption, the
    sampler independently chooses its coins $\delta_{s,a}$ and embeds
    them with the decryption key.  Evaluation of the resulting Boolean
    circuit is deterministic.  Input encryption uses fresh coins at
    posting; no sender erases those coins or its signing state.

    Sender non-commitment lets the simulator post the ciphertext without
    knowing the input and explain the encryption coins when corruption
    reveals it.  The signature remains valid for the unchanged
    ciphertext.  Full NCE supplies both the sender and receiver properties
    we use.

  \item[Zero defaults]
    A missing accepted ciphertext means the all-zero input of the
    prescribed length.  Represent an accepted ciphertext by a public
    flag 1 followed by its fixed-length encoding.  If none is accepted,
    use $\DefaultInput$, a flag 0 followed by a zero payload of the same
    length.  The circuit maps this record directly to $0^{\ell_{s,a}}$.
    It also uses this value if decryption of an accepted ciphertext
    returns $\bot$ or a string of the wrong length.

    This public default record is not a message from the provider.
    Its bits select labels just like ordinary ciphertext bits; no
    correctness assumption for encryption with fixed coins is needed.

  \item[Label committees]
    All inputs to the augmented garbled circuit are now public bit
    strings: the accepted ciphertext records, the incoming state
    ciphertext, and the successor's state public key.  The application
    inputs remain private because they are encrypted inside these records.
    For every garbled input wire $w$, two fresh, disjoint committees
    hold robust secret sharings of $L_{w,0}$ and $L_{w,1}$, respectively.
    The committee-to-bit association is public.
    Each member holds private information for just one committee.

    The committee corresponding to the fixed public bit acts; the other
    remains silent.  The selected committee's honest members publish their
    decryption keys, making their encrypted shares publicly recoverable.
    Robust reconstruction tolerates missing or incorrect shares within
    the threshold.  Honest members of the other committee remain silent;
    its initially corrupt members' shares alone do not reveal its label.
    Later corruption of the speakers does not expose the silent committee.

    If the garbling requires common encoding material $\GarbleAux$
    (Section~\ref{sec:garbling}), bundle a copy with each label before
    sharing it.  The selected committees then release a complete encoded
    input without a separate release step.

  \item[Non-committing delivery and delayed release]
    A committee member's share and public release instructions, or an
    output recipient's pad, are delivered in a receiver non-committing
    ciphertext under its registered key.  Input providers need no such
    private delivery.  Each delivery role receives one bundled
    ciphertext under an unused delivery key in the round's sample and
    decrypts it through the ledger's protected interface.  This is the one
    ciphertext per key programmed using RNCE; other descriptions may
    produce ordinary encryptions under that public key.

    Delivery packets use a fixed public length bound and padding, so
    their keys may be registered before their exact contents or roles
    are known.  On exposure---after speaking for
    honest committee members, or at any time for output recipients---RNCE
    lets the simulator explain an earlier ciphertext, including the
    receiver's key-generation tape.  Together with share completion,
    it lets the proof fix corrupt shares first and complete them to the
    selected label later.

    All input choices for a round are closed before any honest committee
    releases labels and their bundled common material.
    This matches the garbling scheme's separation between an early
    circuit and one later encoding of the complete input.

  \item[Encrypted state between independent garblings]
    The sample for round $s$ generates a fresh public-key encryption pair
    $(\statepk_s,\statesk_s)$.  It publishes $\statepk_s$ and uses
    $\statesk_s$ only inside the garbled circuit, to decrypt the state
    received from the preceding round.  These are state-transfer keys,
    distinct from the registered non-committing delivery keys.

    After computing the new state $\sigma_s$, the round circuit outputs
    \[
      c_s=\Enc(\statepk_{s+1},\sigma_s;\rho_s),
    \]
    where the fresh encryption randomness $\rho_s$ is also hidden inside
    the garbled computation.  The next round decrypts $c_s$ using its own
    embedded secret key.  Thus successive rounds exchange an encryption
    of the state, not correlated wire labels or a secret held by a
    speaking intermediary.  Here we use ordinary IND-CPA encryption.

    Crucially, $\statepk_{s+1}$ and the incoming ciphertext are late
    public inputs, not part of the sampling description: the current round
    can be sampled before either exists.  Committees release only the labels
    for their protocol-specified values.  Publishing both labels would
    allow evaluation with a different successor key.

  \item[Private outputs through public decoding]
    The sampler masks each output bit $y_i$ with an independent uniform
    pad $r_i$ and publishes its RNCE encryption under the output
    recipient's unused delivery key.  The pad is embedded in the round circuit,
    whose publicly decoded output is
    \[
      z_i=y_i\oplus r_i.
    \]
    The recipient privately decrypts $r_i$ and recovers $y_i$; it need not
    post a message, but may still be corrupted at any time.
    The pad encryption is produced during sampling,
    outside the garbled circuit.  This keeps garbled evaluation and
    decoding public while delivering the actual output privately.
\end{description}

\subsection{The construction}
\label{sec:overview-construction}

The protocol prepares each round before evaluating the preceding one.
Preparation fixes a garbled circuit and its private deliveries; evaluation
waits for the round's inputs.  This separation lets the preceding round
encrypt its new state to a successor whose public key is already available,
without requiring either round to know all future circuits or participants.

\begin{description}
  \item[Bootstrap the first nonce]
    The protocol's fixed allocation rule reserves a nonce committee from
    unused registered computation parties.  Everyone derives a nonce-only
    universal sample whose distribution chooses $Q_1$, shares it, and
    encrypts the shares to that committee.  The committee remains silent
    until the first round's other description fields have been fixed.

  \item[Prepare the first round]
    Fix the first application circuit, its input groups and lengths,
    output-wire assignments, registered signing and delivery keys, and
    disjoint label committees for all augmented-circuit input bits.
    Also reserve the committee for $Q_2$.  Use unused registered parties,
    waiting for more registrations only if too few are available.
    Once these fields are frozen on the ledger, release $Q_1$ as described
    above and form $d_1$.  Everyone locally derives the same main sample:
    the garbled circuit, its state-reception public key, its bundle of
    input-encryption public keys, the participants' encrypted information,
    and the encrypted shares of $Q_2$.
    No party has to post this sample.  The first circuit starts with the
    prescribed empty state, rather than decrypting a preceding ciphertext.

  \item[Prepare the successor]
    Before evaluating round $s$, fix the other description fields for
    round $s+1$, including its unused delivery roles and the nonce
    committee for round $s+2$.  The committee already holding shares of
    $Q_{s+1}$ now releases it, and everyone derives $\Sample(U,d_{s+1})$.
    This makes $\statepk_{s+1}$ available
    as a public input to round $s$.  Preparing the successor requires its
    application circuit, public input lengths, and registered role keys,
    but not its private inputs,
    the still-uncomputed ciphertext $c_s$, or $\statepk_{s+2}$.
    There is therefore only one round of lookahead, not a recursive demand
    to prepare the whole future computation.

  \item[Close the current inputs]
    Providers receive their inputs at the round's common start and post
    signed ciphertexts in its input phase; those without inputs stay silent.
    At the deadline, each group's first acceptable ciphertext or
    $\DefaultInput$ fixes its public record.  The incoming state ciphertext
    and successor public key are also fixed by the preceding evaluation
    and successor sample, respectively.  No evaluator chooses these values
    afresh, and all are fixed before honest committees release an encoding.

  \item[Release one encoding]
    Members of all selected committees disclose their decryption keys in
    one synchronous release round.  At its deadline, missing or malformed
    contributions are treated as absent.  Everyone robustly reconstructs
    one label per input wire and its bundled common material.  The
    unselected committees' honest members remain silent permanently;
    later disclosures cannot reopen an input choice.

  \item[Evaluate and deliver]
    Everyone evaluates the garbled circuit locally with this encoding.
    It decrypts the incoming state and accepted application-input
    ciphertexts, substitutes zero strings for default records, applies
    the application circuit, and produces the masked output bits
    and $c_s$, the encryption of the updated state under
    $\statepk_{s+1}$.  Each output recipient privately recovers its pad
    and unmasks its bit.  The masked outputs and $c_s$ are common public
    consequences of the ledger transcript; they need not be posted by
    another party.

  \item[Continue with fresh roles]
    Round $s+1$ is already prepared.  It becomes the current round, and
    the protocol prepares its successor before repeating the input and
    evaluation steps.  Each round uses one ledger-specified description
    and sample.  Unused registered parties supply later roles, with new
    registrations needed only to replenish the pool.  Setup fixes per-round
    circuit and sample-size bounds, but not the total number of rounds.
\end{description}

\paragraph{Input size.}
One provider may supply any string within the round's public length
bound using one ciphertext and one signature.  A constant-rate SNCE
instantiation uses $O(\ell)+\poly(\lambda)$ ciphertext bits for an
$\ell$-bit input.  Thus encrypting a long input need not multiply its
length by a security-parameter factor.  This is a ciphertext-size
statement, not a constant-overhead MPC claim: the public keys, garbled
decryption, and committees have their own costs.  In particular, every
ciphertext bit needs two label committees, of which only one speaks.

\paragraph{When the next application is not yet known.}
The next application circuit or its recipients may depend on the current
result.  In that case, insert an internal buffer round whose application
circuit is the identity on the private state, with no new application
inputs or outputs.  Prepare this buffer and its fresh committee roles
before evaluating the current application.  The current round encrypts
its state to the buffer's public key.  After the result is available,
choose and prepare the next application round; the buffer can then
decrypt and reencrypt the unchanged state to that round's public key.
This costs one extra sample and evaluation, without revealing the state
or adding an application-visible invocation.  Direct lookahead avoids
this overhead whenever the next application is already known.
